\section{Discussion}
\label{sec:discussion}
The strongest evidence concerns the write boundary, not the agent's choice of a setpoint. Across the probed simulator mappings, admissible and inadmissible requests received their expected responses. In the mock-driver fixture, window violations were rejected in the two configurations that retained the window check and accepted in the two that did not, while structural violations were rejected in all four. This supports conformance on the tested branches, with expected responses based on the same configured limits. The endpoint matrix covers writes, not acquisition mappings or parity between task and direct calls.

The distinct records matter when interpreting the exploratory runs: a command outcome is not an independent dispatch record, a simulated readback is not physical actuation, and a later process measurement does not identify the effect of a write. The selected fixed-horizon comparison lacks a documented batch-selection rule and a controller baseline. None of the language-model episodes reached its target, and none retained a correlated acquisition-to-device trace. The runtime defines where an agent can propose and reassess an action; it does not demonstrate autonomous iteration.

A stronger evaluation would first retain a correlated task-mediated trace and compare both entry paths on the same cases, including read-direction checks and post-approval state changes. Process efficacy requires a prespecified study with verified resets, common sampling and horizons, complete trajectories, and a named controller baseline. Retained runtime and model snapshots would make agent trials auditable. Approval effects and physical safety require separate studies.

Two extensions remain future work: autonomous orchestration (a lead agent schedules governed steps and reassesses the results) and process-experience reuse (validated context, intents, and outcomes inform later proposals). Neither would confer write authority or bypass dispatch, approval, or revalidation.

\section{Conclusion}
\label{sec:conclusion}
We presented a configuration-driven runtime for connecting online observations and typed setpoint proposals to scene-specific mappings and governed execution. At 55 writable endpoints, all 220 prespecified admissible requests received the expected acceptance response, all 219 collected structural violations received the expected rejection, and all 216 available simulator readbacks matched. The mock-driver fixture rejected window violations only when the window check was enabled. These tests support conformance of selected write mappings and guards to their configured rules. The scripted comparison remains exploratory, and all three language-model episodes ended below the target band. Neither establishes complete agent-to-device execution, autonomous closed-loop control, comparative efficacy, or physical-plant safety. Separating command outcomes, readbacks, and later process observations makes clear what each record establishes. Autonomous orchestration and process-experience reuse remain future work within this boundary.
